\section*{Capítulo \ref{ch:b2:batteries-electrolysis} --- Baterias, células a combustível e eletrólise}

\begin{solution}{exo:b2:batteries-electrolysis:1}
Zinco sozinho: cerca de \qty{-0.82}{V} e uma corrente de 0.15 (unidades do
modelo). Em contato com o cobre: cerca de \qty{-0.74}{V} e 2.4. A corrente, e
portanto a velocidade de dissolução, é umas dezesseis vezes maior em contato
com o cobre.
\end{solution}

\begin{solution}{exo:b2:batteries-electrolysis:2}
$m = 500 \times 3600 \times 65.38/(2 \times 96\,485) = \qty{610}{g}$.
\end{solution}

\begin{solution}{exo:b2:batteries-electrolysis:3}
$E^\circ = \qty{2.04}{V}$ (como no capítulo); seis células em série dão
\qty{12}{V}.
\end{solution}

\begin{solution}{exo:b2:batteries-electrolysis:4}
$1/6.94 = \qty{0.144}{mol}$; $0.144 \times 96\,485 = \qty{1.39e4}{C}$, isto é,
\qty{3.86}{A.h}.
\end{solution}

\begin{solution}{exo:b2:batteries-electrolysis:5}
A 100~\%: $200 \times 86\,400 \times 63.546/(2 \times 96\,485) = \qty{5.69}{kg}$;
$\eta_F = 5.50/5.69 = 0.967$.
\end{solution}

\begin{solution}{exo:b2:batteries-electrolysis:6}
$w = 2 \times 96\,485 \times 3.3/(0.90 \times 0.06538) = \qty{1.08e7}{J/kg}$, isto é,
\qty{3.0}{kW.h/kg}.
\end{solution}

\begin{solution}{exo:b2:batteries-electrolysis:7}
$U = E - rI$: $r = (1.55 - 1.40)/(0.60 - 0.10) = \qty{0.30}{\ohm}$; $E = 1.55 + 0.30
\times 0.10 = \qty{1.58}{V}$.
\end{solution}

\begin{solution}{exo:b2:batteries-electrolysis:8}
Ânodo \ce{2Cl- -> Cl2 + 2e-}; cátodo \ce{2H2O + 2e- -> H2 + 2OH-}. $\Delta_r G^\circ =
2(-157.24) - 2(-131.23) - 2(-237.13) = \qty{422.2}{kJ/mol}$, $U_{\min} =
\qty{2.19}{V}$. O cloro reagiria com o hidróxido (\ce{Cl2 + 2OH- -> Cl- +
ClO- + H2O}), estragando os dois produtos, e uma mistura de cloro e
hidrogênio pode explodir.
\end{solution}

\begin{solution}{exo:b2:batteries-electrolysis:9}
\ce{NaCl -> Na + 1/2Cl2}, $n = 1$: $U_{\min} = 310\,674/96\,485 = \qty{3.22}{V}$.
Por quilograma de sódio: $96\,485 \times 3.22/0.02299 = \qty{1.35e7}{J}$, isto é,
\qty{3.75}{kW.h}.
\end{solution}

\begin{solution}{exo:b2:batteries-electrolysis:10}
Sobre um cátodo de zinco, a redução de \ce{H+} é lenta e sua barreira fica
abaixo da onda do zinco: à medida que o potencial é abaixado, o zinco se
deposita primeiro, com pouco hidrogênio. Sobre a platina, a redução de \ce{H+}
é rápida e começa perto de \qty{0}{V}: no início só se desprende hidrogênio.
Uma vez a platina recoberta de zinco, o cátodo se comporta como um cátodo de
zinco, e o zinco se deposita.
\end{solution}

\begin{solution}{exo:b2:batteries-electrolysis:11}
A \qty{1250}{K}: $\Delta_f G^\circ(\ce{Al2O3}) = -1328.3 + 0.75 \times 66.0 =
\qty{-1278.8}{kJ/mol}$, $\Delta_f G^\circ(\ce{CO2}) = \qty{-396.1}{kJ/mol}$. $\Delta_r G^\circ =
3(-396.1) - 2(-1278.8) = \qty{1369}{kJ}$, $n = 12$: $U_{\min} = \qty{1.18}{V}$. Ânodo inerte:
$\Delta_r G^\circ = \qty{2557.5}{kJ}$, $U_{\min} = \qty{2.21}{V}$; o carbono que
queima fornece cerca de \qty{1.03}{V} do trabalho.
\end{solution}

\begin{solution}{exo:b2:batteries-electrolysis:12}
Por mol de água, a carga $2F$ é a mesma nos dois sentidos: rendimento $0.70/1.80 =
0.39$. O resto é calor: as \omterm{def:b2:current-potential-curves:overpotential}{sobretensões} dos dois eletrodos de oxigênio (lentos
nos dois sentidos), as dos eletrodos de hidrogênio e as perdas ôhmicas.
\end{solution}

\begin{solution}{pb:b2:batteries-electrolysis:1}
\textbf{1.} Alumínio $+3 \to 0$; o oxigênio continua $-2$; carbono $0 \to +4$.
\textbf{2.} Cátodo \ce{Al^3+ + 3e- -> Al}; ânodo \ce{C + 2O^2- -> CO2 + 4e-}.
\textbf{3.} \ce{2Al2O3 + 3C -> 4Al + 3CO2}, $n = 12$.
\textbf{4.} \ce{Al2O3}: $-1328.3 + 0.75 \times (1328.3 - 1262.3) = \qty{-1278.8}{kJ/mol}$;
\ce{CO2}: \qty{-396.1}{kJ/mol}.
\textbf{5.} $3(-396.1) - 2(-1278.8) = \qty{+1369}{kJ}$.
\textbf{6.} $1\,369\,000/(12 \times 96\,485) = \qty{1.18}{V}$.
\textbf{7.} $2 \times 1278.8/(12 \times 96.485) = \qty{2.21}{V}$: a oxidação do carbono,
ela mesma \omterm{def:b2:reaction-free-energy:exergonic}{exergônica}, paga quase metade do trabalho de decomposição do óxido.
\textbf{8.} $10^6/26.982 = \qty{3.706e4}{mol}$.
\textbf{9.} $3 \times 96\,485 \times 3.706 \times 10^4 = \qty{1.073e10}{C}$.
\textbf{10.} $1.073 \times 10^{10}/0.94 = \qty{1.141e10}{C}$.
\textbf{11.} $1.141 \times 10^{10}/3.0 \times 10^5 = \qty{3.80e4}{s}$, \qty{10.6}{h}.
\textbf{12.} $24/10.6 = 2.27$: \qty{2.27}{t} por dia.
\textbf{13.} Parte do alumínio dissolvido no banho chega ao gás anódico e é
reoxidada pelo dióxido de carbono, de modo que sua carga é gasta duas vezes;
alguma corrente também escapa pelo banho sem eletrolisá-lo.
\textbf{14.} $1.141 \times 10^{10} \times 4.1 = \qty{4.68e10}{J}$, \qty{13.0}{MW.h}.
\textbf{15.} \qty{13.0}{kW.h/kg}.
\textbf{16.} $3 \times 96\,485 \times 1.18/0.026982 = \qty{1.27e7}{J/kg}$, \qty{3.5}{kW.h/kg}.
\textbf{17.} Em calor: as \omterm{def:b2:current-potential-curves:overpotential}{sobretensões} dos eletrodos e, sobretudo, a queda
ôhmica no banho. Esse calor mantém o banho fundido a cerca de
\qty{1250}{K}, o que de outro modo exigiria combustível.
\textbf{18.} $4.1 \times 3.0 \times 10^5 = \qty{1.23}{MW}$.
\textbf{19.} $28\,690/0.026982 = \qty{1.06e6}{J/kg}$, \qty{0.30}{kW.h/kg}: cerca de 2~\% da
energia da \omterm{def:b2:batteries-electrolysis:electrolysis}{eletrólise}.
\textbf{20.} $1.18/4.1 = 0.29$.
\textbf{21.} $3.706 \times 10^4 \times 3/4 \times 12.011 = \qty{334}{kg}$.
\textbf{22.} $3.706 \times 10^4 \times 3/4 \times 44.009 = \qty{1.22e3}{kg}$.
\textbf{23.} O topo quente dos ânodos queima ao ar; parte do carbono sai
como monóxido de carbono (\ce{C + CO2 -> 2CO} a essa temperatura), o que
consome o dobro de carbono por átomo de oxigênio.
\textbf{24.} \qty{1.22}{kg} de \ce{CO2} por quilograma de alumínio.
\textbf{25.} A tensão mínima aumentaria cerca de \qty{1.03}{V}, mas o ânodo
liberaria oxigênio em vez de dióxido de carbono e não seria consumido.
\textbf{26.} A cuba consome $\boldsymbol{\approx \qty{13}{kW.h}}$ de energia
elétrica por quilograma de alumínio.
\end{solution}
